% MT019 / explanation — canonical editable content.
回归预测一个数值。第\(i\)条记录的预测为\(f(x_i)\)，真实值为\(y_i\)，残差\(r_i=f(x_i)-y_i\)就是带正负方向的预测误差；例如预测5、实际3，残差为2，残差平方为4。


普通最小二乘（ordinary least squares，OLS）把各条残差平方后相加，最小化\(\sum_i r_i^2\)。Ridge再加\(\alpha\sum_jw_j^2\)，LASSO再加\(\alpha\sum_j|w_j|\)。二者都控制权重，区别在于压到零附近时的行为。



\begin{itemize}

\item \textbf{A错误。} 更多有代表性的数据通常使估计更稳定，不能说数据越多越容易过拟合；数据来源与模型复杂度仍需考虑。

\item \textbf{B错误。} LASSO恰能使一些系数精确为0，并非只能接近0。

\item \textbf{C正确。} 系数为0时，该特征退出预测式，因此LASSO可用于特征选择。

\item \textbf{D按通常行为正确。} Ridge连续缩小权重，通常不会主动制造精确零。若某特征与目标的相关项本来为0，其系数也可能恰为0，所以“通常”不能省成“永不”。

\item \textbf{E正确。} 令\(\alpha=0\)，两种正则项都消失，目标回到OLS。

\end{itemize}


选读：用一个只有一个权重的教学目标看差别。最小化\(\frac12(w-3)^2+\alpha w^2\)，求导得\((w-3)+2\alpha w=0\)，所以Ridge解为\(w=3/(1+2\alpha)\)，随着惩罚变大逐渐靠近0。


若改成\(\frac12(w-3)^2+\alpha|w|\)，在\(w>0\)处求导得\(w=3-\alpha\)；当\(\alpha\ge3\)，这个正解不再成立，最优值停在尖点\(w=0\)。这就解释了L1的“精确选零”。这个一维例子不是一般多维LASSO都能逐坐标独立求解的保证。


\textbf{出处：} 2021A期中 Q6，原卷第3页（5分）


\textbf{答案依据：} 2021A期中 Q6 随卷／配套参考答案；条件与解释经复核。
